\centerline{\Huge \bf Перекладывание треугольников}

\begin{flushright}
    \scriptsize{Не хотите геому в среду $-$ получайте в воскресенье!}
\end{flushright}

\problem{1} \textit{(Ключевой факт)} В треугольниках $ABC$ и $A'B'C'$ $AB = A'B'$ и $\angle BAC = \angle B'A'C'$ и $BC = B'C'$, что равенство сторон $BC$ и $B'C'$ равносильно равенству углов $\angle BCA$ и $\angle B'C'A'$.

\problem{2} В выпуклом четырехугольнике $ABCD: AD = BC$, $\angle ABD + \angle CDB = 180^\circ$. Докажите, что $\angle BAD = \angle BCD$.

\problem{3} Пусть $K$ — середина стороны $BC$ треугольника $ABC$. На лучах $AB$ и $AC$ взяты точки $X$ и $Y$ соответственно таким образом, что $AX = AY$ и точка $K$ лежит на отрезке $XY$. Докажите, что $BX = CY$.

\problem{4} В четырехугольнике $ABCD$ внешний угол при вершине $A$ равен углу $BCD$, $AD = CD$. Докажите, что $BD$ — биссектриса угла $ABC$.

\problem{5} В выпуклом четырехугольнике $ABCD$, в котором $AB = CD$, на сторонах $AB$ и $CD$ выбраны точки $K$ и $M$ соответственно. Оказалось, что $AM = KC$, $BM = KD$. Докажите, что угол между прямыми $AB$ и $KM$ равен углу между прямыми $KM$ и $CD$.

\problem{6} В неравнобедренном треугольнике $ABC$ биссектрисы $A_1$ и $B_1$ пересекаются в точке $I$. Найдите $I$, если $A_1 = B_1$.

\problem{7} В выпуклом четырехугольнике $ABCD$ $\angle A = \angle C$. На продолжении $BA$ за точкой $A$ выбрана точка $E$ так, что $BC = AE$. Оказалось, что $\angle BDC = \angle ADE$ и $\angle ABD = \angle CAD$. Докажите, что $AC \parallel BD$.

\problem{8} Внутри квадрата $ABCD$ выбрана точка $P$, а на сторонах $BC$ и $CD$ — точки $K$ и $L$ соответственно таким образом, что $PK = PL$. Точка $Q$ отрезка $BK$ такова, что $BQ = DQ$, а углы $BKP$ и $DPL$ равны. Докажите, что $PQ$ перпендикулярна $DP$.

\problem{9} На сторонах $BC$ и $AB$ треугольника $ABC$ выбраны точки $A'$ и $C'$ соответственно. При этом $A'A$ и $C'C$ пересекаются в точке $K$. Оказалось, что $A'C' = CA'$ и $\angle ABC = \angle A'KC$. На отрезке $KC'$ выбрана точка $P$ такая, что $2PK = A'K + KC'$. Докажите, что $\angle APC = 90^\circ$.

\problem{10} Пусть $O$ — точка пересечения диагоналей $C$ и $AD$ выпуклого пятиугольника $ABCDE$. Известно, что $AB = DE$, $BC = AD$ и $\dfrac{CE}{2} = \dfrac{CA}{2}$. Докажите, что $AO = OE$.

